% !TeX program = xelatex
% !TeX root = main.tex
\section{贡献与评价}
\sectionlead{section-5}{五、贡献与评价}{浏览轨迹数据，给资产定价提供了哪些认识？}

\subsection{5.1 证据与局限}
\begin{frame}[label=scorecard]{五项研究结论怎样约束模型的信息输入}
  \crumb{Part 5}{5.1 证据与局限 · 记分卡}{论文 I、VIII · 汇报人综合}
  \begin{columns}[T,onlytextwidth]
    \column{.34\textwidth}\footnotesize
    \textbf{五项观测事实}\par\smallskip
    \leadrow{\jump{depth}{研究较薄}}{5.7 分钟，50 户交易标的中位数／笔。}
    \leadrow{\jump{topics}{常看价格}}{价格图覆盖 93\%，484 户。}
    \leadrow{\jump{windows}{默认窗口}}{约 73\% 的访问是 1 日图。}
    \leadrow{\jump{pc-two}{风格有别}}{快／慢信息与标的倾向同变。}
    \leadrow{\jump{timing}{临近下单}}{同股研究在交易附近增加。}
    \column{.64\textwidth}\footnotesize
    \setlength{\tabcolsep}{3pt}
    \renewcommand{\arraystretch}{1.06}
    \begin{tabularx}{\linewidth}{@{}L{.29\linewidth}Y@{}}
      \toprule 文献输入 & 观测证据与判断\\\midrule
      MV, CAPM & 显式风险输入少见\\
      DSSW & 薄且偏价格，部分一致\\
      Hong--Stein & 风格有别，新闻和价格同侧\\
      BSV & 部分一致：盈利覆盖 40\%\\
      收益外推 & 价格输入一致\\
      股利信号 & 覆盖 17\%\\
      limited attention & 财报、事件有关注峰值\\
      同质输入 & 投入右偏\\
      \bottomrule
    \end{tabularx}
  \end{columns}
\end{frame}

\subsection{5.2 研究评价}
\begin{frame}[label=evaluation]{研究评价}
  \crumb{Part 5}{5.2 研究评价}{汇报人评价}
  \leadrow{数据与测度}{利用跨域名 URL，刻画投资者的信息获取行为。}
  \bigskip
  \leadrow{理论联系}{用微观数据的注意力、主题和时点事实，对资产定价模型的信息输入假设作经验比较。}
  \bigskip
  \leadrow{不足}{未联接完整持仓、信念与实际成交，难以识别信息获取对决策的因果影响。}
\end{frame}

\newcommand{\riskcell}[2]{\cellcolor{wine!#1}#2}
\begin{frame}[label=limitations]{汇报总结}
  \crumb{Part 5}{5.2 研究评价 · 汇报总结}{论文 II--VIII · 汇报综合}
  {\small\bfseries\color{slate}个人投资者实际看什么、研究多久，又怎样接近交易？}\par\smallskip
  \small
  \setlength{\tabcolsep}{4pt}
  \renewcommand{\arraystretch}{1.02}
  \begin{tabularx}{\linewidth}{@{}>{\bfseries\color{wine}}L{.18\linewidth}Y@{}}
    \toprule\rowcolor{wine!10} 研究环节 & \textbf{本文的回答}\\\midrule
    数据与方法 & 用 \textbf{2007 年 484 户}的跨站 URL，识别研究股票、内容、用时及订单时点。\\
    \rowcolor{wine!5}总体画像 & 研究投入较薄，价格信息最常见：高匹配 \textbf{50 户}交易标的的中位数为 \textcolor{wine}{\textbf{5.7 分钟／笔}}；\textbf{484 户}中 \textcolor{wine}{\textbf{93\%}} 看过价格图。\\
    研究差异 & 大股票被研究得更多；股票特征对主题倾斜的解释较弱。投资者之间存在跨主题参与和快／慢信息的风格差异。\\
    \rowcolor{wine!5}交易时点 & 同股研究在临近下单时集中增加，多数研究路径涉及的主题较少。\\
    研究贡献 & 直接观察个人投资者的信息获取行为，为资产定价模型的信息集假设提供微观证据。\\\bottomrule
  \end{tabularx}
\end{frame}

\begin{frame}[label=thanks]
  \vfill\centering
  {\LARGE\bfseries\color{wine}谢谢大家}\par\bigskip
  {\large 请各位老师和同学批评指正}\par\bigskip
  {\small\jump{limitations}{返回 · 汇报总结}}
  \vfill
\end{frame}
